\begin{proof}
First of all, since $\Gamma$ is undirected and connected, we have, by Lemmas \ref{lem:connectivity} and~\ref{lem:undirected}, that $\rho$ is a self-dual, faithful representation of $G$. Next, the condition that $\Gamma$ is a tree implies that $$\abs{\{\text{undirected edges in } \Gamma\}} = \abs{\mathrm{Irr}(G)}-1 = \abs{G//G} -1.$$

By Corollary \ref{cor:trace} we have:
 \begin{equation}\label{eq:sq_char}
\dim \Hom_G(\rho^{\otimes 2}, \C[G]^{adj}) = tr(A_{\Gamma}^{\otimes 2})= 2(\abs{G//G} -1).  
 \end{equation}

 Now, 
 $$\C[G]^{adj} = \bigoplus_{[g] \in G//G} \Span\{g |g\in [g]\} \cong \bigoplus_{[g] \in G//G} \C[G/C(g)] \cong \bigoplus_{[g] \in G//G} \C\uparrow^G_{C(g)}, $$
 where $C(g)$ is the centralizer of an element $g$ in the conjugacy class $[g]$ (the isomorphisms depend on the choices of representatives $g\in [g]$, of course).
 
 
 Thus we have:
\begin{align}\label{eq:sum_dim_end}
 \nonumber&2(\abs{G//G} -1) = \dim \Hom_G(\rho^{\otimes 2}, \C[G]^{adj})  = \sum_{[g] \in G//G} \dim \Hom_G(\rho^{\otimes 2}, \C\uparrow^G_{C(g)} )  \\ &=\sum_{[g] \in G//G} \dim \Hom_{C(g)}\left(\rho^{\otimes 2} \downarrow_{C(g)}, \triv\right)= \sum_{[g] \in G//G}  \dim \End_{C(g)}(\rho \downarrow_{C(g)})
\end{align}
(the last equality follows from the fact that $\rho$ is self-dual). 


Let us compute $ \dim \End_{C(g)}(\rho \downarrow_{C(g)})$. This value is clearly at least $1$, and we will show that this is precisely $2$ when $g\notin Z(G)$.


Assume $\rho$ is reducible as a $G$-representation. Hence it is reducible as a $C(g)$-representation for each $g\in G$. This implies: $$\dim \End_{C(g)}(\rho \downarrow_{C(g)})\geq 2$$ for all $g\in G$ and so $$\sum_{[g] \in G//G}  \dim \End_{C(g)}(\rho \downarrow_{C(g)}) \geq 2\abs{G//G} >2(\abs{G//G} -1)$$ which is a contradiction. Hence $\rho$ is irreducible.  

For $g\in Z(G)$ we have: $C(g)=G$ so clearly $ \dim \End_{C(g)}(\rho \downarrow_{C(g)})=1$ in that case.


Now let $g\in G \setminus Z(G)$. The operator $\rho(g)$ is an intertwining operator on the $C(g)$-representation $\rho \downarrow_{C(g)}$ and hence acts by scalar on any simple $C(g)$-summand of $\rho \downarrow_{C(g)}$. Note that $\rho(g)$ cannot be a scalar endomorphism: indeed, if that were the case, $\rho(g)$ would commute with $\rho(h)$ for any $g\in G$. Since $\rho$ is faithful, this would imply that $g \in Z(G)$, contrary to our assumption on $g$.




Hence $\rho(g)$ is not a scalar endomorphism, and so $\rho \downarrow_{C(g)}$ is not simple.
 
Hence we have: for any $g\in G \setminus Z(G)$, $\dim \End_{C(g)}(\rho \downarrow_{C(g)}) \geq 2$. 
 
 
 Next, by Corollary \ref{cor:bipartite_center_Z_2} we have: $Z(G) = C_2$. So by Equation \eqref{eq:sum_dim_end}, we have:
 $$\sum_{\substack{[g] \in G//G,\; g \notin Z(G)}} \dim \End_{C(g)}(\rho \downarrow_{C(g)}) = 2(\abs{G//G} -2).$$
 
 Notice that the left hand side has $(\abs{G//G} -2)$ summands, so that the average value of $\dim \End_{C(g)}(\rho \downarrow_{C(g)}) $ is $2$. But since we showed that the minimal value is also $2$, we conclude that $\dim \End_{C(g)}(\rho \downarrow_{C(g)}) =2$ for any $g\notin Z(G)$, i.e. the representation $\rho \downarrow_{C(g)}$ has length~$2$.
 
 If $C(g)$ is abelian for some $g\notin Z(G)$, then $\rho \downarrow_{C(g)}$ is a direct sum of two $1$-dimensional representations and the statement of the theorem is proved. If $G$ itself is abelian, then $\dim \rho = 1$ and $\Gamma$ is tree only if $G= C_2$, $\rho = \sgn$; again, the statement of the theorem is then proved. 
 
Thus we will assume from now on that $C(g)$ is not abelian for any $g\in G$.
 
Since $\dim \End_{C(g)}(\rho \downarrow_{C(g)}) =2$, for any $g\notin Z(G)$, $\rho(g)$ is a diagonalizable endomorphism with precisely $2$ distinct eigenvalues. Furthermore, since $\rho \cong \rho^*$ as $G$-representations, we have: for each eigenvalue $z$ of $\rho(g)$, its complex conjugate $\bar{z}$ is also an eigenvalue of $\rho(g)$ with the same multiplicity. 

So the set of eigenvalues of each $\rho(g), g\notin Z(G)$ is either $\{\pm 1\}$ or $\{z, \bar{z}\}$ where $z$ is a root of unity of order $\abs{G}$, $z\neq \pm 1$. Note that the second option can occur only if $\dim \rho$ is even.

Assume {that the eigenvalues of $\rho(g)$ belong to $\{\pm 1\}$ for all $g\in G$}. This implies that each {(non-trivial)} element in $G$ has order $2$. In particular, $$\forall g, h\in G, \, ghg^{-1}h^{-1} = ghgh=1_G \; \Rightarrow \; gh=hg$$ so $G$ is abelian. But that contradicts our assumption. 

So there exists $g \in G$ such that $g^2 \neq 1_G$, i.e. the eigenvalues of $\rho(g)$ are $z_1 \neq z_2$ (these are roots of unity, and $z_1 =\bar{z_2} { \notin \{\pm 1\}}$). Let $V_{z_k}$ denote the eigenspace of $\rho(g)$ corresponding to the eigenvalue $z_k$, $k=1,2$; then $V_{z_1} \oplus V_{z_2}$ is the entire space underlying the representation $\rho$, and $\dim V_{z_1} = \dim V_{z_2}$. 

Let $h\in C(g)$. The operators $\rho(h), \rho(g)$ commute so they are simultaneously diagonalizable ($\rho(h)$ also having at most $2$ distinct eigenvalues). We will say that $h$ is {\it aligned} with $g$ if the restriction of $\rho(h)$ to each of the subspaces $V_{z_k}, k=1,2$ is a scalar endomorphism. The set of elements $h\in C(g)$ aligned with $g$ clearly forms an abelian subgroup.

Since we assumed that the group $C(g)$ is itself not abelian, we can find $h \in C(g)$ which is not aligned with $g$. We have: $h\notin Z(G)$, so we denote by $t_1 \neq t_2$ the distinct eigenvalues of $\rho(h)$. Since $h$ is not aligned with $g$, the eigenvalues of $\rho(gh)$ would be $\{t_j z_k |k, j\in \{1,2\}\}$ or at least $3$ of these products. {Since $\rho(gh)$ also has at most $2$ distinct eigenvalues, this} implies that in at least one of the pairs $(t_1 z_1, t_2 z_2)$, $(t_1 z_2, t_2 z_1)$ the elements coincide. {Without loss of generality, let us assume that $t_1z_1 = t_2z_2 = \bar{t_1}\bar{z_1}$. So the eigenvalues of $\rho(gh)$ are either $\{t_1z_1, t_1z_2\}$ or $\{t_1z_1, t_2z_1\}$. Then $$ (t_1 z_1)^2= t_1 z_1 t_2 z_2 =\abs{t_1}^2 \abs{z_1}^2 = 1 .$$ 
So $t_1z_1 = \pm 1$ (implying $t_1 = \pm z_2$) and the eigenvalues of $\rho(gh)$ must be $\pm 1$. Thus at least one of the following equalities holds: $ z_2^2 = \pm t_1 z_2 = \pm 1$ or $z_1^2 = \pm t_2 z_1 = \pm 1$, any of which implies} $z_1, z_2 \in \{\pm i\}$. So we see that for every $g\in G$, the eigenvalues of $\rho(g)$ are either $\{\pm i\}$ or a subset of $\{\pm 1\}$.

In particular, $\rho(g)^2 = \pm \id$ for each $g\in G$, so $G/Z(G)$ is a group {whose elements all square to identity}. {Hence $G/Z(G)$ is} of the form $\mathbb{F}_2^m$ (an elementary abelian $2$-group). Since $Z(G)=C_2$, we conclude that $G$ is an extra special $2$-group, $m$ is even and $\rho$ is the unique irreducible representation of $G$ of dimension greater than $1$ (see Example~\ref{ex:extra_special}). This concludes the proof of the theorem.

\end{proof}
